Die exakte Zunahme der Gravitationsenergie beim Heben um $h$ ist
\begin{align}
 \Delta E &= G\,M_\mathrm{E}\,m\left(\frac{1}{r_\mathrm{E}}-\frac{1}{r_\mathrm{E}+h}\right)
   = G\,\frac{M_\mathrm{E}\,m}{r_\mathrm{E}}\cdot\frac{h}{r_\mathrm{E}+h}
\end{align}
Mit der Fallbeschleunigung an der Erdoberfläche $g = G\,M_\mathrm{E}/r_\mathrm{E}^2$ wird daraus
\begin{align}
 \Delta E &= m\,g\,h\cdot\frac{r_\mathrm{E}}{r_\mathrm{E}+h}
\end{align}
Die Näherung $m\,g\,h$ ist also zu gross; ihr relativer Fehler ist
\begin{align}
 p &= \frac{m\,g\,h-\Delta E}{\Delta E} = \frac{r_\mathrm{E}+h}{r_\mathrm{E}}-1 = \frac{h}{r_\mathrm{E}}
\end{align}
Für $p = \pFeO$ und $r_\mathrm{E} = \RErO$ (FoTaBe S.~209):
\begin{align}
 h &= \hFF = \pFeO\cdot\REr = \hF \approx \result{\hFP{3}{2}}
\end{align}
